% ------------------------ CONCLUSION ---------------------------------
\section{Conclusion}
\label{sec:conc}

OSEBX rebalancing is predictable and partly tradable before the official announcement date. Using only public Euronext rule-book variables, GLM and XGBoost models achieve AUROC above 0.97 across the 2010--2023 out-of-sample window. More importantly, portfolios formed from predicted additions and deletions generate economically meaningful alpha after transaction costs, short-borrow costs and factor adjustment.

The strongest lesson is that model selection must follow the investment objective. The naive persistence model is the most accurate classifier but creates no portfolio. The conditional-threshold XGBoost model is less accurate, yet produces more diversified event exposure and is the only specification that remains significant inside a realistic 2\% tracking-error enhanced-index sleeve. For a portfolio manager, that is the relevant victory: a repeatable signal that can be scaled inside a benchmarked risk budget.

The result is bounded by three caveats. First, the Norwegian factor data end in May 2022, so the latest rebalancing events are included in classification and portfolio plots but not in the FF3 regressions. Second, the cost assumptions are institutional; retail implementation would face wider spreads, higher commissions and less reliable short access. Third, the index effect can decay as arbitrage capital learns the rule. The natural extension is therefore a live paper-trading implementation over the next five OSEBX rebalances, with pre-registered thresholds and execution rules. Subject to those caveats, the strategy is implementable today with public index methodology, standard market data and an open-source machine-learning workflow.
